\chapter{Электрические основы: КМОП, выходы и шины}
\label{ch:electrical}

До сих пор мы считали, что логический элемент --- идеальный чёрный ящик: на входе
0 или 1, на выходе 0 или 1. Для разработки на ПЛИС этого почти достаточно, но на
границе кристалла с внешним миром идеализация перестаёт работать. Ножка ПЛИС
подключается к светодиодам и кнопкам, к чужим микросхемам, к длинным проводам.
Здесь важно знать, из чего сделаны логические элементы, как они «толкают» ток
и что значит неопределённый уровень. Эта глава --- короткий мост между
булевой алгеброй и реальными проводами.

\section{Транзистор как выключатель}

Современные цифровые схемы строят из полевых транзисторов типа \term{МОП}
(металл---оксид---полупроводник, MOSFET). Для нашей цели транзистор можно
представить управляемым ключом с тремя выводами: затвор (управляющий),
исток и сток (между ними течёт ток).

\begin{itemize}
  \item \textbf{n-канальный} транзистор (nMOS) проводит ток, когда на затворе
        \emph{высокий} уровень (1), и разрывает цепь при низком.
  \item \textbf{p-канальный} транзистор (pMOS) работает наоборот: он
        проводит при \emph{низком} уровне на затворе и закрыт при высоком.
\end{itemize}

\section{КМОП-инвертор}

Если соединить оба транзистора, как показано на рис.~\ref{fig:cmos_inv}, получится
инвертор. pMOS подключён к питанию $U_{\text{ПИТ}}$, nMOS --- к земле. Затворы
соединены вместе и образуют вход, а соединённые стоки образуют выход.

\begin{figure}[H]
\centering
\begin{tikzpicture}[font=\sffamily\small, sw/.style={draw, minimum width=11mm, minimum height=9mm, fill=white, rounded corners=1pt}]
  \draw (0,3.4) -- (0,3.8) node[above] {$U_{\text{ПИТ}}$};
  \node[sw] (p) at (0,2.7) {pMOS};
  \node[sw] (n) at (0,0.9) {nMOS};
  \draw (0,3.4) -- (p.north);
  \draw (p.south) -- (0,1.8) -- (n.north);
  \draw (0,1.8) -- (1.4,1.8) node[right] {выход};
  \draw (n.south) -- (0,0) node[below] {земля};
  \draw (-1.4,1.8) node[left] {вход} -- (-0.85,1.8) -- (-0.85,2.7) -- (p.west);
  \draw (-0.85,1.8) -- (-0.85,0.9) -- (n.west);
  \fill (-0.85,1.8) circle (1.2pt);
  \fill (0,1.8) circle (1.2pt);
  % таблица
  \node[align=left, anchor=west] at (3.5,3.6) {Вход = 1: nMOS открыт, pMOS закрыт\\{}$\Rightarrow$ выход соединён с землёй (0)};
  \node[align=left, anchor=west] at (3.5,2.2) {Вход = 0: pMOS открыт, nMOS закрыт\\{}$\Rightarrow$ выход соединён с питанием (1)};
  \node[align=left, anchor=west, text=FpgaGreen!50!black] at (3.5,0.6) {В любом устойчивом состоянии открыт\\ровно один транзистор: тока от питания\\к земле нет.};
\end{tikzpicture}
\caption{КМОП-инвертор: пара транзисторов, один из которых всегда закрыт}
\label{fig:cmos_inv}
\end{figure}

Главное достоинство \term{КМОП} (комплементарных МОП, CMOS) в том, что в
покое схема почти не потребляет энергию: нет пути от питания к земле. Ток
появляется только в момент переключения, когда нужно перезарядить
ёмкость нагрузки и входы следующих элементов. Поэтому расход энергии
пропорционален частоте.

\begin{ideabox}
Мощность, которую расходует КМОП-схема на переключение,
$P \approx C\cdot U^2\cdot f$, где $C$ --- суммарная ёмкость переключаемых
цепей, $U$ --- напряжение питания, $f$ --- частота переключений. Чтобы
уменьшить мощность, надо снижать напряжение (оно стоит в квадрате!),
частоту или число переключающихся узлов. Именно поэтому в ПЛИС
хорошая привычка --- не тактировать то, что не работает.
\end{ideabox}

\section{Элементы И-НЕ и ИЛИ-НЕ}

Чтобы построить элемент с двумя входами, добавляют транзисторы
по принципу двойственности. В элементе \term{И-НЕ} два nMOS стоят
последовательно (цепь к земле замыкается, только если оба входа равны 1), а
два pMOS --- параллельно (достаточно одного нуля, чтобы выход соединился с
питанием). У элемента \term{ИЛИ-НЕ} всё наоборот: nMOS параллельно, а pMOS
последовательно.

\begin{figure}[H]
\centering
\begin{tikzpicture}[font=\sffamily\small, sw/.style={draw, minimum width=9mm, minimum height=8mm, fill=white, rounded corners=1pt}]
  % И-НЕ
  \node at (1.5,4.9) {\textbf{И-НЕ}};
  \draw (1.5,4.5) -- (1.5,4.2) node[above, yshift=-2pt] {};
  \draw (0.4,4.2) -- (2.6,4.2);
  \node[sw] (p1) at (0.4,3.4) {p};
  \node[sw] (p2) at (2.6,3.4) {p};
  \draw (0.4,4.2) -- (p1.north); \draw (2.6,4.2) -- (p2.north);
  \draw (0.4,3.0) -- (0.4,2.7) -- (2.6,2.7) -- (2.6,3.0);
  \draw (1.5,2.7) -- (3.4,2.7) node[right] {$Y$};
  \node[sw] (n1) at (1.5,1.9) {n};
  \node[sw] (n2) at (1.5,0.7) {n};
  \draw (1.5,2.7) -- (n1.north); \draw (n1.south) -- (n2.north); \draw (n2.south) -- (1.5,0.2) node[below] {земля};
  \node[left] at (-0.3,3.4) {$A$}; \node[left] at (0.9,1.9) {$A$}; \node[right] at (3.2,3.4) {$B$}; \node[left] at (0.9,0.7) {$B$};
  % ИЛИ-НЕ
  \begin{scope}[xshift=7.5cm]
  \node at (1.5,4.9) {\textbf{ИЛИ-НЕ}};
  \node[sw] (q1) at (1.5,4.0) {p};
  \node[sw] (q2) at (1.5,2.8) {p};
  \draw (1.5,4.6) -- (q1.north); \draw (q1.south) -- (q2.north);
  \draw (1.5,2.4) -- (1.5,2.0);
  \draw (1.5,2.0) -- (3.4,2.0) node[right] {$Y$};
  \draw (0.4,2.0) -- (2.6,2.0);
  \node[sw] (m1) at (0.4,1.2) {n};
  \node[sw] (m2) at (2.6,1.2) {n};
  \draw (0.4,2.0) -- (m1.north); \draw (2.6,2.0) -- (m2.north);
  \draw (0.4,0.8) -- (0.4,0.3) -- (2.6,0.3) -- (2.6,0.8);
  \draw (1.5,0.3) -- (1.5,0.0) node[below] {земля};
  \node[left] at (0.9,4.0) {$A$}; \node[left] at (0.9,2.8) {$B$}; \node[left] at (-0.0,1.2) {$A$}; \node[right] at (3.1,1.2) {$B$};
  \end{scope}
\end{tikzpicture}
\caption{Схемы КМОП-элементов. Показаны только транзисторы и соединения; затворы с одинаковой буквой соединены между собой. Питание подключено к верхнему выводу}
\end{figure}

Поэтому в КМОП «базовые» элементы --- инвертирующие. Для элемента И
нужно добавить инвертор (всего шесть транзисторов вместо четырёх). Именно по
этой причине в главе~\ref{ch:gates} мы называли И-НЕ и ИЛИ-НЕ универсальными и
экономными: они ещё и физически проще.

\section{Уровни сигналов и запас помехоустойчивости}

Реальный выход не выдаёт «ровно ноль» и «ровно питание»: он
гарантирует лишь, что напряжение 0 будет не выше $U_{OL}$, а 1 --- не ниже
$U_{OH}$. Вход, в свою очередь, считает низким всё, что не выше $U_{IL}$, и
высоким всё, что не ниже $U_{IH}$. В промежутке находится запрещённая зона.

\begin{table}[H]
\centering
\caption{Типичные пороги для логики 3,3~В (LVTTL/LVCMOS33)}
\small
\begin{tabular}{llr}
\toprule
\textbf{Параметр} & \textbf{Смысл} & \textbf{Значение} \\
\midrule
$U_{OL}$ & выход при 0, не выше & 0,4~В \\
$U_{OH}$ & выход при 1, не ниже & 2,4~В \\
$U_{IL}$ & вход читает 0, если не выше & 0,8~В \\
$U_{IH}$ & вход читает 1, если не ниже & 2,0~В \\
\bottomrule
\end{tabular}
\end{table}

Запас на помехи --- это разность между тем, что гарантирует выход, и тем,
что требует вход: $N_L = U_{IL}-U_{OL} = 0{,}4$~В для нуля и
$N_H = U_{OH}-U_{IH} = 0{,}4$~В для единицы. Наводка меньше этих значений
не испортит сигнал. Именно поэтому цифровая схема так устойчива: пока помехи
меньше запаса, сигнал «очищается» на каждом элементе.

\section{Нагрузка и задержка}

К выходу элемента подключены входы следующих элементов, у каждого из которых
есть небольшая ёмкость, плюс сам провод. Чем больше ёмкость, тем дольше выход
переходит из 0 в 1, и тем больше задержка. Число входов, которые может
запитать один выход, называют \term{коэффициентом разветвления}. В ПЛИС этот
эффект видно в отчёте о частоте: сигнал, идущий на сотню входов (например,
сброс), медленнее сигнала, идущего на один.

Для ножек микросхемы основное ограничение --- ток. Светодиод требует нескольких
миллиампер, и в файле ограничений мы задаём \code{DRIVE=8} (до 8~мА). Подключать
к ножке мощную нагрузку (мотор, реле) нельзя: нужен усилитель из транзистора.

\section{Типы выходов}

В цифровой технике встречаются три разных схемы выхода.

\begin{figure}[H]
\centering
\begin{tikzpicture}[scale=0.82, transform shape, font=\sffamily\small, sw/.style={draw, minimum width=8mm, minimum height=7mm, fill=white, rounded corners=1pt}]
  % Push-pull
  \node at (0,4.6) {\textbf{Двухтактный}};
  \draw (0,3.2) -- (0,3.0);
  \node[sw] (a) at (0,2.5) {p}; \node[sw] (b) at (0,1.3) {n};
  \draw (0,3.2) node[above] {$U_{\text{ПИТ}}$} -- (a.north); \draw (a.south) -- (b.north); \draw (0,1.9) -- (1.3,1.9) node[right] {выход}; \draw (b.south) -- (0,0.7) node[below] {};
  \node[align=center] at (0,0.2) {сам выдаёт и 0, и 1};
  % Open drain
  \begin{scope}[xshift=5.2cm]
  \node at (0,4.6) {\textbf{Открытый сток}};
  \node[sw] (c) at (0,1.3) {n};
  \draw (0,1.9) -- (c.north); \draw (0,1.9) -- (1.3,1.9) node[right] {выход}; \draw (c.south) -- (0,0.7);
  \draw (0,1.9) -- (0,2.4) node[sw, anchor=south] {$R$} ; \draw (0,3.1) -- (0,3.3) node[above] {$U_{\text{ПИТ}}$};
  \node[align=center] at (0,0.2) {тянет только к нулю,\\единицу даёт резистор};
  \end{scope}
  % Three-state
  \begin{scope}[xshift=10.6cm]
  \node at (0,4.6) {\textbf{Три состояния}};
  \node[draw, minimum width=18mm, minimum height=14mm, fill=white] (t) at (0,1.9) {буфер};
  \draw (-1.6,1.9) node[left] {вход} -- (t.west);
  \draw (t.east) -- (1.5,1.9) node[right] {выход};
  \draw (0,3.2) node[above] {$OE$} -- (t.north);
  \node[align=center] at (0,0.2) {0, 1 или «отключён» (Z)};
  \end{scope}
\end{tikzpicture}
\caption{Три схемы выходов цифровых элементов}
\end{figure}

\begin{itemize}
  \item \textbf{Двухтактный} (push-pull) выход активно подаёт и 0, и 1. Такие
        выходы стоят на большинстве ножек ПЛИС.
  \item \textbf{С открытым стоком} (open drain) умеет только соединять выход с
        землёй. Единицу создаёт внешний подтягивающий резистор. Зато к одной
        линии можно подключить сколько угодно таких выходов: линия будет в
        нуле, если \emph{хоть один} притянул её к земле, то есть получается
        функция И (<<монтажное И>>). Так устроена шина I$^2$C (глава~ef{ch:uart}).
  \item \textbf{С тремя состояниями} (tri-state, 3-state) может выдавать 0, 1
        или быть \emph{отключённым}: высокоимпедансное состояние $Z$. Именно
        отключаемый выход позволяет нескольким источникам по очереди
        пользоваться общим проводом.
\end{itemize}

\section{Шина с несколькими источниками}

\term{Шина} --- это группа общих проводов, к которым подключены несколько
устройств. В каждый момент передаёт ровно одно, остальные должны быть
отключены (состояние $Z$). Выбор передатчика --- та же задача, что решает
мультиплексор (глава~\ref{ch:mux}), только физическая: вместо логического
выбора выход всех, кроме одного, просто отсоединяется от провода.

\begin{figure}[H]
\centering
\begin{tikzpicture}[font=\sffamily\small]
  \foreach \i/\y in {1/2.4, 2/1.2, 3/0} {
    \node[draw, minimum width=14mm, minimum height=7mm, fill=white] (b\i) at (2,\y) {буфер $\i$};
    \draw (0.1,\y) node[left] {$D_\i$} -- (b\i.west);
    \draw (b\i.east) -- (4.2,\y);
    \node[font=\sffamily\scriptsize] at (2,\y+0.58) {$OE_\i$};
  }
  \draw[very thick] (4.2,-0.4) -- (4.2,2.8);
  \draw[very thick] (4.2,1.2) -- (5.5,1.2) node[right] {общая шина};
  \node[align=left, anchor=west] at (7.0,2.0) {Разрешён только один из\\$OE_1, OE_2, OE_3$.};
  \node[align=left, anchor=west, text=FpgaRed] at (7.0,0.3) {Если включить два с разными\\данными, выходы «борются»:\\ток короткого замыкания.};
\end{tikzpicture}
\caption{Три источника на одной шине}
\end{figure}

\begin{warnbox}
Если два передатчика одновременно включены и выдают разные значения, один тянет
линию к питанию, другой к земле: через них течёт большой ток, а на линии
получается неопределённое напряжение. Это \term{конфликт на шине}. В
симуляции он выглядит как значение \code{X}, на настоящей плате --- как
перегрев и выход из строя. Управляющая схема обязана гарантировать, что
одновременно включён один передатчик.
\end{warnbox}

\section{Выход с тремя состояниями на практике}

В языках описания аппаратуры такой выход записывают условным присваиванием:
«если выход разрешён, подать значение, иначе --- состояние $Z$». Двунаправленная
ножка добавляет к этому чтение того, что фактически стоит на проводе, чтобы
схема могла и передавать, и принимать по одной и той же линии. Готовый пример
с описанием, схемой и проверкой мы разберём в главе~\ref{ch:bridge}.

\begin{infobox}[Что бывает внутри ПЛИС]
Внутри кристалла нет общих проводов с несколькими передатчиками: каждый
сигнал питается ровно одним источником, а выбор источника выполняют
мультиплексоры. Элементы с тремя состояниями есть только в блоках ввода-вывода
на ножках (примитивы \code{TBUF} и \code{IOBUF}). Поэтому описывать внутри схемы
общие шины с состоянием $Z$ не надо: синтезатор их всё равно заменит
мультиплексорами.
\end{infobox}

\section{Подтяжки}

Вход, который ни к чему не подключён, «плавает»: его уровень определяется
наводками, и схема может работать случайно. Поэтому к каждому входу нужен
резистор, который задаёт уровень по умолчанию: \term{подтяжка к питанию}
(pull-up) или \term{к земле} (pull-down). В ПЛИС такие резисторы
встроены, и включаются они в файле \code{.cst} параметром \code{PULL\_MODE}.
Для кнопки с замыканием на землю подтяжка нужна к питанию (\code{UP}), а для
кнопки с замыканием на питание --- к земле (\code{DOWN}). Мы это делали в
проекте~1 главы~\ref{ch:examples}.

\section{Вопросы и задачи}
\label{sec:tasks-electrical}

Ответы собраны в приложении~\ref{app:answers}.

\begin{taskbox}[Задачи к главе]
\begin{enumerate}
  \item Нарисуйте КМОП-схему элемента ИЛИ-НЕ с тремя входами: сколько
        транзисторов каждого типа потребуется и как они соединены?
  \item Схема потребляет $P = C U^2 f$. Как изменится мощность, если
        снизить питание с 3,3~В до 1,8~В, оставив частоту прежней?
  \item Вычислите мощность переключения цепи ёмкостью 10~пФ при питании
        3,3~В и частоте 27~МГц.
  \item Выход выдаёт напряжение 2,6~В, а вход требует не менее 2,0~В. Каков запас
        помехоустойчивости для единицы?
  \item Почему на шине I$^2$C используют выходы с открытым стоком и подтягивающие
        резисторы, а не двухтактные выходы?
  \item Что произойдёт, если два буфера с тремя состояниями одновременно
        выдадут на общую линию 0 и 1?
\end{enumerate}
\end{taskbox}
